% !TEX program = pdflatex
\documentclass[12pt,a4paper]{article}

% ── Packages ──────────────────────────────────────────────────────────────────
\usepackage[utf8]{inputenc}
\usepackage[T1]{fontenc}
\usepackage{amsmath,amssymb,amsfonts}
\usepackage{mathtools}
\usepackage{booktabs}
\usepackage{graphicx}
\usepackage{hyperref}
\usepackage[margin=1in]{geometry}
\usepackage{natbib}
\usepackage{multirow}
\usepackage{caption}
\usepackage{subcaption}
\usepackage{xcolor}
\usepackage{url}
\usepackage{threeparttable}
\usepackage{algorithm}
\usepackage{algorithmic}
\usepackage{array}

\hypersetup{
  colorlinks=true,
  linkcolor=blue!70!black,
  citecolor=blue!70!black,
  urlcolor=blue!70!black,
}

% ── Title ─────────────────────────────────────────────────────────────────────
\title{%
  Optimal Linear Prediction from Fractional Gaussian Noise:\\
  A Split-Architecture Forecaster that Beats ARIMA%
}

\author{
  Rick Galbo\\
  Wayy Research\\
  \texttt{rick@wayyresearch.com}
}

\date{\today}

% ══════════════════════════════════════════════════════════════════════════════
\begin{document}
\maketitle

% ── Abstract ──────────────────────────────────────────────────────────────────
\begin{abstract}
We present a split-architecture forecaster that separates point prediction
from uncertainty quantification using the autocovariance structure of
fractional Gaussian noise (fGn).  The point forecast is the Wiener--Hopf
best linear unbiased predictor (BLUP), which solves
$\boldsymbol{\Gamma}\mathbf{a} = \boldsymbol{\gamma}_h$ for optimal
weights on past returns, where $\boldsymbol{\Gamma}$ is the fGn
autocovariance matrix and $\boldsymbol{\gamma}_h$ is the cross-covariance
at horizon~$h$.  Uncertainty intervals are probability-weighted quantiles
of Monte Carlo paths generated by a novel simulation pipeline:
pattern-conditioned seeding from historical analogues, a numerically
stable integrated fBm kernel, and fractal consistency scoring that
weights paths by Hurst match, volatility match, and pattern continuity.
Controlled experiments on synthetic fractional Brownian motion with known
Hurst exponents $H \in \{0.3, 0.4, 0.5, 0.6, 0.7\}$ show that the
split forecaster (v3) achieves lower RMSE than ARIMA at \emph{all}
tested~$H$ (ratios 0.96--0.98), higher directional accuracy at all~$H$
(52.5--69.2\% vs.\ 15.8--66.7\%), higher Sharpe ratio at $H \geq 0.4$,
and well-calibrated 95\% coverage (90.8--96.7\%).  We document two
intermediate architectures (v1: symmetric path averaging; v2: directional
tilt) and show why each fails, providing a constructive path from
distributional forecasting to optimal linear prediction.
\end{abstract}

\smallskip
\noindent\textbf{Keywords:} Fractional Gaussian noise $\cdot$
Wiener--Hopf prediction $\cdot$ Hurst exponent $\cdot$
Pattern-conditioned Monte Carlo $\cdot$
Fractal consistency scoring $\cdot$ ARIMA

% ══════════════════════════════════════════════════════════════════════════════
\section{Introduction}
\label{sec:introduction}

The Hurst exponent~$H$ characterises the long-range dependence of a time
series~\citep{Hurst1951,MandelbrotVanNess1968}.  When $H \neq 0.5$, the
increments of fractional Brownian motion (fBm) exhibit non-zero
autocorrelation: positive for $H > 0.5$ (persistence) and negative for
$H < 0.5$ (anti-persistence).  Financial time series frequently exhibit
$H \neq 0.5$~\citep{Peters1994,Cont2001}, motivating fractal-based
forecasting.

The \textsc{FracTime} library~\citep{Galbo2026} implements a
probability-weighted Monte Carlo forecaster with a novel simulation
architecture: paths are seeded by historical pattern matching (not
independent fBm draws), scored by fractal consistency against the
observed Hurst exponent and volatility, and aggregated via
probability-weighted quantiles.  Walk-forward experiments on real market
data revealed a characteristic pattern: FracTime produced higher RMSE than
ARIMA but better directional accuracy and better-calibrated prediction
intervals.

This paper investigates \emph{why} this trade-off arises and \emph{how} to
resolve it.  We show that:
\begin{enumerate}
  \item The original architecture (v1) averages symmetric Monte Carlo paths,
        producing near-neutral point forecasts that discard the directional
        information encoded in~$H$.
  \item A directional tilt (v2) partially recovers directional signal but
        still suffers from the fundamental limitation of path averaging.
  \item A split architecture (v3) that uses the \emph{Wiener--Hopf optimal
        linear predictor} for the point forecast and \emph{Monte Carlo
        quantiles} for uncertainty intervals resolves the trade-off,
        beating ARIMA on RMSE, direction, and Sharpe simultaneously.
\end{enumerate}

We validate these findings on controlled synthetic data with known Hurst
exponents, isolating the effect of~$H$ from confounding factors present in
real data.

% ══════════════════════════════════════════════════════════════════════════════
\section{Background}
\label{sec:background}

\subsection{Fractional Brownian Motion and Gaussian Noise}

Fractional Brownian motion $B^H(t)$ is the unique Gaussian process with
stationary increments, zero mean, and covariance:
\begin{equation}
  \operatorname{Cov}\bigl(B^H(t), B^H(s)\bigr) =
    \tfrac{1}{2}\bigl(|t|^{2H} + |s|^{2H} - |t-s|^{2H}\bigr)
  \label{eq:fbm_cov}
\end{equation}
The increments $X_k = B^H(k) - B^H(k-1)$ form \emph{fractional Gaussian
noise} (fGn), a stationary process with autocovariance:
\begin{equation}
  \gamma(k) = \tfrac{\sigma^2}{2}\bigl(|k+1|^{2H} - 2|k|^{2H}
              + |k-1|^{2H}\bigr)
  \label{eq:fgn_acov}
\end{equation}
where $\sigma^2 = \operatorname{Var}(X_k)$.  For $H > 0.5$, $\gamma(k) > 0$
for all $k \geq 1$ and decays as $k^{2H-2}$, producing long memory.  For
$H < 0.5$, $\gamma(1) < 0$ and $|\gamma(k)|$ decays rapidly, concentrating
the anti-persistent signal at short lags.  The Hurst exponent can be
estimated via rescaled range (R/S) analysis~\citep{Hurst1951} or
detrended fluctuation analysis (DFA)~\citep{Peng1994}, with multifractal
extensions~\citep{Kantelhardt2002} for time-varying~$H$.

\subsection{The ARIMA Baseline}

Auto-ARIMA~\citep{BoxJenkinsReinsel2015} selects the order $(p,d,q)$
that minimises AIC and produces the conditional mean forecast:
\begin{equation}
  \hat{P}_{t+h}^{\text{ARIMA}} =
    \mathbb{E}[P_{t+h} \mid P_1, \ldots, P_t]
  \label{eq:arima_forecast}
\end{equation}
This is the minimum-variance linear predictor under the ARIMA model
assumptions, with prediction intervals based on Gaussian innovations.

\subsection{The FracTime Simulation Architecture}
\label{sec:fractime_arch}

FracTime's forecaster generates distributional forecasts through a
three-stage pipeline that differs from standard Monte Carlo in several
respects.

\paragraph{Stage 1: Pattern-conditioned path generation.}
Rather than simulating independent fBm paths from scratch, FracTime first
searches the historical return series for \emph{analogues}---subsequences
whose recent shape resembles current market conditions.  Let
$\mathbf{r}_{t-\ell:t} = (r_{t-\ell+1}, \ldots, r_t)$ be the most recent
$\ell = 21$ returns.  The system scans the history for windows
$\mathbf{r}_{s:s+\ell}$ with normalised correlation above a threshold
$\rho_{\min} = 0.3$:
\begin{equation}
  \hat{\rho}(s) =
    \frac{1}{\ell}\sum_{k=1}^{\ell}
    \frac{(r_{s+k} - \bar{r}_s)}{\hat{\sigma}_s}
    \cdot
    \frac{(r_{t-\ell+k} - \bar{r}_t)}{\hat{\sigma}_t}
  \label{eq:pattern_sim}
\end{equation}
Each match yields a \emph{continuation}: the returns that historically
followed the matched pattern.  Path~$i$ is seeded by sampling a match
(weighted by similarity) and replaying its continuation with additive
Gaussian noise scaled to half the current volatility:
\begin{equation}
  r^{(i)}_{\tau} = r_{s+\ell+\tau}^{\text{hist}} +
    \varepsilon_\tau, \quad
    \varepsilon_\tau \sim \mathcal{N}(0,\, \hat{\sigma}/2)
  \label{eq:pattern_path}
\end{equation}
When no matches exceed~$\rho_{\min}$, the forecaster falls back to pure
fBm path generation using the integrated kernel described below.

\paragraph{Stage 2: Fractal consistency scoring.}
Each generated path receives a probability weight based on how well its
\emph{statistical shape} matches the observed fractal structure.  The score
combines three components via a weighted sum:
\begin{equation}
  w_i = \alpha_H\, f_H(S_i)
       + \alpha_\sigma\, f_\sigma(S_i)
       + \alpha_c\, f_c(S_i)
  \label{eq:fractal_scoring}
\end{equation}
where:
\begin{itemize}
  \item $f_H(S_i) = (1 + 5|\hat{H}(S_i) - H_{\text{target}}|)^{-1}$
        penalises paths whose R/S Hurst deviates from the estimated~$H$;
  \item $f_\sigma(S_i) = (1 + 2|\log(\hat{\sigma}_i / \sigma_{\text{target}})|)^{-1}$
        penalises volatility mismatch; and
  \item $f_c(S_i) = (1 + 10\operatorname{std}(\Delta^2 r^{(i)}))^{-1}$
        rewards path smoothness (pattern continuity) via the second
        difference of returns.
\end{itemize}
Default weights are $\alpha_H = 0.3$, $\alpha_\sigma = 0.3$,
$\alpha_c = 0.4$, emphasising pattern continuity.  Weights are normalised
to a probability distribution: $p_i = w_i / \sum_j w_j$.  This creates a
\emph{non-uniform ensemble} where paths that preserve the observed fractal
geometry are amplified.

\paragraph{Stage 3: Probability-weighted quantiles.}
Prediction intervals are computed as \emph{weighted quantiles} of the
scored path ensemble.  For target quantile~$q$, paths at each step~$t$
are sorted by terminal price, and the quantile is found where the
cumulative probability mass reaches~$q$:
\begin{equation}
  Q_q^{(t)} = S_{(\pi_k)}(t) \quad \text{where} \quad
    \sum_{j=1}^{k} p_{\pi_j} \geq q \cdot \sum_{j=1}^{N} p_{\pi_j}
  \label{eq:weighted_quantile}
\end{equation}
with $\pi$ the permutation that sorts paths by price at step~$t$.
This yields intervals that respect the non-uniform weighting, unlike
standard equally-weighted quantiles.

\paragraph{Integrated fBm kernel.}
When generating fBm paths (the fallback), FracTime uses a numerically
stable variant of the Mandelbrot--Van Ness kernel.  The standard kernel
computes weights $w(i,j) = (i-j)^{H-0.5} - (i-j-1)^{H-0.5}$, which
produces the singularity $0^{H-0.5} = 0^{\text{negative}}$ for $H < 0.5$.
FracTime instead uses the \emph{integrated kernel}:
\begin{equation}
  w(i,j) = (i-j)^{H+0.5} - (i-j-1)^{H+0.5}
  \label{eq:integrated_kernel}
\end{equation}
Since $H + 0.5 > 0$ for all valid~$H$, the term $0^{H+0.5} = 0$ and no
singularity arises.  For $H = 0.5$ all weights equal~$1$, recovering
standard Brownian motion.  For $H > 0.5$, recent noise is weighted
more heavily (persistence); for $H < 0.5$, distant noise retains weight
(anti-persistence in the increments).

\subsection{Directional Information Content}

Raw directional accuracy can be misleading.  A model predicting the correct
direction 39.7\% of the time contains \emph{more} directional information
than one at 57.5\%, because $|39.7 - 50| = 10.3\% > |57.5 - 50| = 7.5\%$.
The 39.7\% model is simply inverted: flipping its signals yields 60.3\%
accuracy.

We therefore define the \emph{directional information content}:
\begin{equation}
  I_{\text{dir}} = \bigl|\text{Dir.\ Acc.} - 0.5\bigr|
  \label{eq:dir_info}
\end{equation}
and require that forecasters achieve Dir.\ Acc.\ $> 0.5$ to be
directionally useful without signal inversion.

% ══════════════════════════════════════════════════════════════════════════════
\section{The Problem: Why Path Averaging Fails}
\label{sec:problem}

\subsection{Architecture v1: Symmetric Path Averaging}

The original FracTime forecaster uses the full pipeline described in
Section~\ref{sec:fractime_arch}: pattern-conditioned path generation,
fractal consistency scoring~(Eq.~\ref{eq:fractal_scoring}), and
probability-weighted quantile intervals.  The point forecast is the
weighted mean of the scored ensemble:
\begin{equation}
  \hat{P}^{\text{v1}}_{t+h} =
    \frac{\sum_{i=1}^N p_i \cdot S_i(h)}{\sum_{i=1}^N p_i}
  \label{eq:v1_forecast}
\end{equation}
The scoring function~$w_i$ (Eq.~\ref{eq:fractal_scoring}) depends on
path \emph{shape}---Hurst match, volatility match, smoothness---but is
\emph{symmetric with respect to direction}: an upward path and its
mirror-image downward path receive identical weights because all three
scoring components ($f_H$, $f_\sigma$, $f_c$) are invariant to the sign
of returns.  Consequently, the weighted mean converges to a value near
the current price $P_t$, yielding a near-neutral forecast regardless of
the true Hurst exponent.

This explains the original observation: FracTime's point forecast has
high RMSE (it predicts approximately $P_t$, the random-walk forecast)
while its \emph{intervals} are well-calibrated---the pattern-conditioned
paths with fractal scoring do capture the correct dispersion, and
the weighted quantile intervals (Eq.~\ref{eq:weighted_quantile}) adapt
to asymmetry and fat tails in the path distribution.

\subsection{Architecture v2: Directional Tilt}

To inject directional signal, v2 introduces a momentum-persistence
interaction.  Let $\bar{r}$ be the mean of recent log-returns, $\hat{\sigma}$
the recent volatility, and $p = 2H - 1$ the persistence parameter.  The
directional signal is:
\begin{equation}
  d = p \cdot \frac{\bar{r}}{\hat{\sigma}}
  \label{eq:v2_signal}
\end{equation}
Paths whose terminal direction agrees with $\operatorname{sgn}(d)$ receive
higher weight:
\begin{equation}
  w_i^{\text{dir}} = \begin{cases}
    1 + \tanh(|d|) & \text{if } \operatorname{sgn}(S_i(h) - P_t) =
                      \operatorname{sgn}(d) \\
    1 - 0.5\tanh(|d|) & \text{otherwise}
  \end{cases}
  \label{eq:v2_weights}
\end{equation}

V2 replaces the symmetric scoring~(Eq.~\ref{eq:fractal_scoring}) with a
four-component score that adds directional alignment ($w_{\text{dir}} = 0.4$)
while retaining Hurst match, volatility match, and smoothness at reduced
weight ($0.2$ each).  This modestly improves Sharpe at high $H$ but remains
fundamentally limited: the weighted mean of a biased-but-asymmetric path
distribution is still pulled toward the centre.  Reweighting paths cannot
produce a point forecast as sharp as a direct prediction from the
autocovariance structure.

% ══════════════════════════════════════════════════════════════════════════════
\section{The Solution: Split Architecture (v3)}
\label{sec:solution}

\subsection{Design Principle}

We decouple point prediction from uncertainty quantification:
\begin{itemize}
  \item \textbf{Point forecast}: Wiener--Hopf BLUP from the fGn
        autocovariance.
  \item \textbf{Prediction intervals}: Empirical quantiles of
        directionally-weighted Monte Carlo paths.
\end{itemize}
This exploits the theoretical optimality of the BLUP for point forecasts
while retaining the distribution-free calibration of MC intervals.

\subsection{Wiener--Hopf Optimal Linear Prediction}
\label{sec:wiener_hopf}

Given $p$ observed returns $r_t, r_{t-1}, \ldots, r_{t-p+1}$, the BLUP
for the return $h$ steps ahead is:
\begin{equation}
  \hat{r}_{t+h} = \sum_{j=0}^{p-1} a_j^{(h)}\, r_{t-j}
  \label{eq:blup}
\end{equation}
where the weight vector $\mathbf{a}^{(h)}$ solves the Wiener--Hopf
equation:
\begin{equation}
  \boldsymbol{\Gamma}\,\mathbf{a}^{(h)} = \boldsymbol{\gamma}_h
  \label{eq:wiener_hopf}
\end{equation}
with $\Gamma_{ij} = \gamma(|i-j|)$ the autocovariance matrix and
$[\boldsymbol{\gamma}_h]_j = \gamma(h+j)$ the cross-covariance vector.
Both are computed from the exact fGn autocovariance
(Eq.~\ref{eq:fgn_acov}) using the estimated Hurst exponent.

The multi-step price forecast accumulates predicted returns:
\begin{equation}
  \hat{P}_{t+h} = P_t \exp\!\Biggl(\sum_{k=1}^{h} \hat{r}_{t+k}\Biggr)
  \label{eq:price_forecast}
\end{equation}

\subsection{Adaptive Lookback}
\label{sec:adaptive_lookback}

The number of past returns $p$ used in the BLUP should match the effective
memory length of the process.  We set $p$ to the largest lag $k$ where the
autocovariance exceeds a threshold:
\begin{equation}
  p = \max\bigl\{k : |\gamma(k)| \geq \epsilon\bigr\},
  \quad \epsilon = 0.01\sigma^2
  \label{eq:adaptive_lookback}
\end{equation}
with a minimum of $p = 3$ for numerical stability.  This yields:

\begin{table}[htbp]
  \centering
  \caption{Adaptive Lookback by Hurst Exponent}
  \label{tab:lookback}
  \begin{tabular}{ccrl}
    \toprule
    $H$ & $\gamma(1)/\sigma^2$ & $p$ & Character \\
    \midrule
    0.3 & $-0.242$ & 6  & Strong reversal at lag~1, fast decay \\
    0.4 & $-0.129$ & 6  & Mild reversal, fast decay \\
    0.5 & $\phantom{-}0.000$ & 3  & No signal (random walk) \\
    0.6 & $+0.149$ & 23 & Mild persistence, slow decay \\
    0.7 & $+0.320$ & 64 & Strong persistence, very slow decay \\
    \bottomrule
  \end{tabular}
\end{table}

This structure captures a key asymmetry: anti-persistent processes
concentrate their predictive signal at lag~1, while persistent processes
benefit from many lags of slowly-decaying memory.

\subsection{Probability-Weighted Monte Carlo Intervals}

Prediction intervals use the same weighted quantile mechanism as v1/v2
(Eq.~\ref{eq:weighted_quantile}), with the directionally-tilted v2
scoring applied to the path ensemble:
\begin{equation}
  \text{CI}_{1-\alpha} = \bigl[Q^{(h)}_{\alpha/2},\;
                                Q^{(h)}_{1-\alpha/2}\bigr]
  \label{eq:mc_ci}
\end{equation}
where $Q^{(h)}_q$ is the probability-weighted quantile at horizon~$h$.
Unlike ARIMA's Gaussian intervals, these are distribution-free and adapt
to skewness, fat tails, and multimodality.  The fractal consistency
scoring ensures that paths contributing to the interval boundaries
preserve the observed Hurst exponent and volatility structure.

% ══════════════════════════════════════════════════════════════════════════════
\section{Experimental Design}
\label{sec:design}

\subsection{Synthetic Data Generation}

We generate fBm price series using the integrated Mandelbrot--Van Ness
kernel~(Eq.~\ref{eq:integrated_kernel}), which is numerically stable for
all $H \in (0,1)$---critical for this study's anti-persistent ($H < 0.5$)
test cases.  The fBm path is differenced to obtain fGn increments, which
are normalised and scaled to the target daily volatility:
\begin{equation}
  r_t = \frac{B^H_t - B^H_{t-1}}{\operatorname{std}(\Delta B^H)}
        \cdot \frac{\sigma_a}{\sqrt{252}}
  \label{eq:synthetic_returns}
\end{equation}
with $\sigma_a = 20\%$ (annualised) and zero drift.

\subsection{Walk-Forward Protocol}

\begin{table}[htbp]
  \centering
  \caption{Experiment Parameters}
  \label{tab:params}
  \begin{tabular}{ll}
    \toprule
    \textbf{Parameter} & \textbf{Value} \\
    \midrule
    Series length ($n$)      & 500 \\
    Training window           & 252 days (expanding) \\
    Step size                 & 21 days \\
    Forecast horizon          & 5 days \\
    Monte Carlo paths         & 500 \\
    Seeds per $H$             & 10 \\
    Hurst values              & $\{0.3, 0.4, 0.5, 0.6, 0.7\}$ \\
    Models                    & v1, v2, v3, ARIMA (auto) \\
    \bottomrule
  \end{tabular}
\end{table}

Each seed generates a different realisation of fBm prices.  Metrics are
averaged across seeds using the sample mean and Bessel-corrected standard
deviation.  The total experiment comprises $5 \times 10 \times 4 = 200$
walk-forward runs with approximately 12 forecast points each, yielding
$\sim$120 forecasts per model per~$H$.

\subsection{Evaluation Metrics}

\begin{itemize}
  \item \textbf{RMSE}: Root mean squared error of point forecasts.
  \item \textbf{Directional Accuracy}: Fraction where
        $\operatorname{sgn}(\hat{P} - P_t)
        = \operatorname{sgn}(P_{t+h} - P_t)$.
  \item \textbf{Directional Information}: $|$Dir.\ Acc.\ $- 0.5|$.
  \item \textbf{Sharpe Ratio}: Annualised Sharpe of a simple long/short
        strategy: $\text{signal}_t = \operatorname{sgn}(\hat{P} - P_t)$,
        $\text{return}_t = \text{signal}_t \cdot (P_{t+h} - P_t) / P_t$,
        annualised by $\sqrt{12}$ (forecasts occur every 21 trading days,
        giving $\sim$12 observations per year).
  \item \textbf{Coverage-95\%}: Fraction of actuals within the 95\%
        prediction interval, following~\citet{GneitingRaftery2007}.
\end{itemize}

% ══════════════════════════════════════════════════════════════════════════════
\section{Results}
\label{sec:results}

\subsection{Main Comparison: v3 vs.\ ARIMA}

Table~\ref{tab:main_results} presents the primary results averaged over
10 seeds per Hurst value.

\begin{table}[htbp]
  \centering
  \caption{v3 (Wiener--Hopf) vs.\ ARIMA: 10 Seeds per~$H$, 5-Day Horizon}
  \label{tab:main_results}
  \begin{threeparttable}
  \begin{tabular}{c cc cc cc cc}
    \toprule
    & \multicolumn{2}{c}{\textbf{RMSE}} &
      \multicolumn{2}{c}{\textbf{Dir.\ Acc.}} &
      \multicolumn{2}{c}{\textbf{Sharpe}} &
      \multicolumn{2}{c}{\textbf{Cov-95\%}} \\
    \cmidrule(lr){2-3} \cmidrule(lr){4-5} \cmidrule(lr){6-7} \cmidrule(lr){8-9}
    $H$ & v3 & ARIMA & v3 & ARIMA & v3 & ARIMA & v3 & ARIMA \\
    \midrule
    0.3 & \textbf{1.87} & 1.91 & \textbf{52.5\%} & 50.8\% &
          0.21 & \textbf{0.36} & \textbf{96.7\%} & 94.2\% \\
    0.4 & \textbf{2.20} & 2.24 & \textbf{55.0\%} & 40.0\% &
          \textbf{0.36} & $-0.08$ & \textbf{95.8\%} & 95.0\% \\
    0.5 & \textbf{2.61} & 2.73 & \textbf{56.7\%} & 15.8\% &
          \textbf{0.71} & $-0.12$ & 92.5\% & \textbf{91.7\%} \\
    0.6 & \textbf{3.10} & 3.22 & \textbf{61.7\%} & 60.0\% &
          \textbf{1.17} & 1.11 & \textbf{92.5\%} & 90.8\% \\
    0.7 & \textbf{3.67} & 3.80 & \textbf{69.2\%} & 66.7\% &
          \textbf{1.94} & 1.84 & 90.8\% & \textbf{91.7\%} \\
    \bottomrule
  \end{tabular}
  \begin{tablenotes}\small
    \item Best per row in bold.  RMSE ratio (v3/ARIMA) ranges from 0.959 to
          0.983.  All direction accuracies for v3 are above 50\% (no signal
          inversion needed).
  \end{tablenotes}
  \end{threeparttable}
\end{table}

\paragraph{RMSE.}
V3 achieves lower RMSE than ARIMA at \emph{every} Hurst value, with ratios
ranging from 0.959 ($H = 0.5$) to 0.983 ($H = 0.4$).  The improvement is
largest at $H = 0.5$ and $H = 0.7$, where the BLUP's adaptive lookback
correctly identifies zero signal (minimum prediction) and deep memory
(many-lag prediction), respectively.

\paragraph{Directional accuracy.}
V3 achieves Dir.\ Acc.\ above 50\% at all $H$ values, meaning its
directional signal is usable without inversion.  The most striking result
is at $H = 0.5$ (pure random walk), where v3 achieves 56.7\% while ARIMA
achieves only 15.8\%.  ARIMA's 15.8\% represents high directional
information ($|15.8 - 50| = 34.2\%$) but in the \emph{wrong} direction,
suggesting systematic misidentification of the process as mean-reverting.

\paragraph{Sharpe ratio.}
V3 dominates ARIMA at $H \geq 0.4$.  The Sharpe advantage grows with $H$:
at $H = 0.7$ (strong persistence), v3 achieves 1.94 versus ARIMA's 1.84.
ARIMA wins only at $H = 0.3$ (Sharpe 0.36 vs.\ 0.21), where the
anti-persistent signal is genuinely weak ($\gamma(1) = -0.24$, fast decay).
Sharpe ratios are annualised with $\sqrt{12}$ (monthly forecast frequency).

\paragraph{Coverage.}
Both models achieve reasonable coverage at all $H$ values.  V3's MC
intervals are slightly better calibrated at low $H$ (96.7\% at $H = 0.3$
vs.\ target 95\%) and comparable at high $H$.

\subsection{Evolution Across Architectures}

Table~\ref{tab:evolution} shows how each scoring version improved upon
its predecessor.

\begin{table}[htbp]
  \centering
  \caption{Architecture Evolution: RMSE and Direction Across Versions}
  \label{tab:evolution}
  \begin{threeparttable}
  \begin{tabular}{c cccc cccc}
    \toprule
    & \multicolumn{4}{c}{\textbf{RMSE}} &
      \multicolumn{4}{c}{\textbf{Dir.\ Acc.}} \\
    \cmidrule(lr){2-5} \cmidrule(lr){6-9}
    $H$ & v1 & v2 & v3 & ARIMA & v1 & v2 & v3 & ARIMA \\
    \midrule
    0.3 & 1.97 & 1.98 & \textbf{1.87} & 1.91
        & 51.7\% & 50.8\% & \textbf{52.5\%} & 50.8\% \\
    0.4 & 2.38 & 2.38 & \textbf{2.20} & 2.24
        & 47.5\% & 46.7\% & \textbf{55.0\%} & 40.0\% \\
    0.5 & 2.87 & 2.88 & \textbf{2.61} & 2.73
        & 44.2\% & 45.8\% & \textbf{56.7\%} & 15.8\% \\
    0.6 & 3.44 & 3.43 & \textbf{3.10} & 3.22
        & 47.5\% & 50.0\% & \textbf{61.7\%} & 60.0\% \\
    0.7 & 4.20 & 4.16 & \textbf{3.67} & 3.80
        & 51.7\% & 57.5\% & \textbf{69.2\%} & 66.7\% \\
    \bottomrule
  \end{tabular}
  \begin{tablenotes}\small
    \item V1 and v2 direction accuracies frequently fall below 50\%,
          requiring signal inversion to be useful.  V3 is above 50\% at
          all~$H$.
  \end{tablenotes}
  \end{threeparttable}
\end{table}

Key observations:
\begin{enumerate}
  \item \textbf{V1 $\to$ v2}: Minimal RMSE change; modest directional
        improvement at $H \geq 0.6$ (from 47.5\% to 50.0\% at $H = 0.6$,
        from 51.7\% to 57.5\% at $H = 0.7$).  Path reweighting alone
        cannot overcome the symmetric-averaging problem.
  \item \textbf{V2 $\to$ v3}: Large RMSE improvement (7--12\% reduction)
        \emph{and} large directional improvement.  The split architecture
        resolves the RMSE--direction trade-off by using the right tool for
        each objective.
\end{enumerate}

\subsection{Directional Information Content}

Table~\ref{tab:dir_info} reports $I_{\text{dir}} = |\text{Dir.\ Acc.} - 0.5|$,
the true measure of directional signal strength.

\begin{table}[htbp]
  \centering
  \caption{Directional Information Content $|\text{Dir.\ Acc.} - 50\%|$}
  \label{tab:dir_info}
  \begin{tabular}{c cccc}
    \toprule
    $H$ & v1 & v2 & v3 & ARIMA \\
    \midrule
    0.3 &  8.3\% & 10.8\% &  5.8\% &  5.8\% \\
    0.4 & 10.8\% & 13.3\% &  6.7\% & 11.7\% \\
    0.5 & 12.5\% & 14.2\% & 10.0\% & \textbf{34.2\%} \\
    0.6 & 12.5\% & 11.7\% & \textbf{11.7\%} & 15.0\% \\
    0.7 & 10.0\% &  9.2\% & \textbf{19.2\%} & 18.3\% \\
    \bottomrule
  \end{tabular}
\end{table}

At $H = 0.5$ (random walk), ARIMA has the highest directional information
(34.2\%) but at 15.8\% accuracy---the signal is inverted.  V3's 10.0\% at
56.7\% is genuinely predictive without inversion.  At $H = 0.7$
(strong persistence), v3 leads with 19.2\% information, all above 50\%.

This distinction is critical: a model with 40\% directional accuracy and
10\% information content requires an external mechanism to detect and
correct the inversion.  V3's consistent accuracy above 50\% means its
signals are directly actionable.

% ══════════════════════════════════════════════════════════════════════════════
\section{Analysis}
\label{sec:analysis}

\subsection{Why the BLUP Beats Conditional Mean ARIMA}

The BLUP~(Eq.~\ref{eq:blup}) is optimal for fGn by construction.  ARIMA
must estimate the order $(p,d,q)$ from data, introducing model selection
error.  On fBm-generated prices, three factors favour the BLUP:

\begin{enumerate}
  \item \textbf{Correct model class.}  fGn is an ARFIMA(0,$d$,0) process
        with $d = H - 0.5$.  Auto-ARIMA fits integer-order models to a
        fractionally-integrated process, incurring systematic
        misspecification.
  \item \textbf{Known autocovariance.}  The BLUP uses the exact fGn
        autocovariance parameterised by the estimated~$H$.  ARIMA must
        estimate all AR/MA coefficients from the training window,
        introducing sampling error.
  \item \textbf{Adaptive memory length.}  The autocovariance-driven
        lookback~(Eq.~\ref{eq:adaptive_lookback}) matches the effective
        memory of the process.  ARIMA's order selection via AIC may under-
        or over-fit the memory structure.
\end{enumerate}

\subsection{The Anti-Persistent Challenge}

At $H = 0.3$, v3 achieves only 52.5\% directional accuracy (information
content 5.8\%), barely above chance.  This reflects a fundamental
difficulty: the anti-persistent signal is concentrated at lag~1
($\gamma(1) = -0.24$) but the signal-to-noise ratio per forecast is low.
Over a 5-day horizon, the reversal at lag~1 is partially offset by weaker
reversals at lags 2--5, diluting the cumulative signal.

ARIMA faces the same challenge: 50.8\% accuracy at $H = 0.3$.  Neither
model extracts strong directional signal from weakly anti-persistent data.
This is consistent with the theoretical prediction interval: the BLUP
prediction variance at $H = 0.3$ is only marginally below the unconditional
variance.

\subsection{Separation of Concerns}

The split architecture's key insight is that \emph{point prediction} and
\emph{uncertainty quantification} have different optimal solutions:

\begin{itemize}
  \item The BLUP minimises MSE by exploiting the exact autocovariance
        structure.  It is sharp (low RMSE) but provides no distributional
        information.
  \item The FracTime pipeline (pattern-conditioned paths, fractal scoring,
        weighted quantiles) provides rich distributional
        information---quantiles, skewness, tail behaviour---but the weighted
        mean of even a well-scored ensemble is a poor point forecast because
        the scoring criteria (Hurst match, volatility match, smoothness)
        are symmetric with respect to direction.
  \item By combining the BLUP for the point forecast with
        probability-weighted MC quantiles for intervals, v3 gets the best
        of both approaches: an MSE-optimal centre and distribution-free
        tails.
\end{itemize}

This is analogous to the distinction between the mean and quantile
functions in quantile regression: each serves a different loss function,
and conflating them degrades both.

% ══════════════════════════════════════════════════════════════════════════════
\section{Discussion}
\label{sec:discussion}

\subsection{Practical Implications}

The results suggest a clear model selection framework:

\begin{enumerate}
  \item \textbf{When $H$ is estimable and $\neq 0.5$}: Use v3
        (Wiener--Hopf BLUP + MC intervals).  The BLUP exploits the
        autocovariance directly, and adaptive lookback ensures
        appropriate memory length.
  \item \textbf{When $H \approx 0.5$}: The BLUP correctly predicts
        near-zero returns (no false signal), while ARIMA may overfit
        short-term noise.  V3 still provides well-calibrated intervals.
  \item \textbf{For risk management}: MC intervals from all versions
        outperform ARIMA's Gaussian intervals, particularly at low~$H$
        where v3 achieves 96.7\% coverage versus ARIMA's 94.2\%.
\end{enumerate}

\subsection{Limitations}

\begin{enumerate}
  \item \textbf{Hurst estimation error.}  The BLUP assumes the true~$H$
        is known.  In practice, $H$ is estimated from historical data with
        uncertainty.  Misestimation of~$H$ degrades the BLUP toward a
        random-walk forecast.
  \item \textbf{fGn assumption.}  Real financial returns are not exactly
        fGn.  They exhibit fat tails, volatility clustering, and regime
        changes that violate the Gaussian, stationary-increment
        assumptions.  The BLUP is optimal only within the fGn model class.
  \item \textbf{Synthetic validation.}  These results are on controlled
        fBm data where the BLUP's model is correctly specified.  Real-data
        performance will depend on how well $H$ captures the actual
        dependence structure.
  \item \textbf{Small sample sizes.}  With 10 seeds and $\sim$12 forecasts
        per seed, the standard error on directional accuracy is
        approximately 4.6\%.  Some observed differences (e.g., 1.7\% margin
        at $H = 0.3$) are not statistically significant.
        Diebold--Mariano tests~\citep{DieboldMariano1995} on the RMSE
        differences would strengthen the claims but require larger samples
        for adequate power.
\end{enumerate}

\subsection{Future Work}

Several extensions are natural:

\begin{enumerate}
  \item \textbf{Real-data validation} across equities, crypto, forex,
        and commodities using the walk-forward framework in
        \texttt{fractime.research}.
  \item \textbf{Regime-switching Hurst} where $H$ varies over time.
        The BLUP can be re-computed at each forecast step using the
        rolling Hurst estimate.
  \item \textbf{Fat-tailed innovations} where the Gaussian assumption
        of fGn is replaced by Student-$t$ marginals while preserving
        the autocovariance structure.
  \item \textbf{Multi-step recursive prediction} where intermediate
        predicted returns are fed back into the BLUP for longer horizons.
  \item \textbf{Ensemble with ARIMA} combining the BLUP point forecast
        (for fractal signal) with ARIMA (for short-memory signal) via
        inverse-variance weighting.
\end{enumerate}

% ══════════════════════════════════════════════════════════════════════════════
\section{Reproducibility}
\label{sec:reproducibility}

All experiments are fully reproducible:
\begin{verbatim}
  pip install fractime==0.6.0
  python scripts/run_scoring_experiment.py --seeds 10
\end{verbatim}
Results are saved to
\texttt{paper\_results/scoring\_experiment/results.json}.

\noindent Source code:
\url{https://github.com/Wayy-Research/fracTime}

% ══════════════════════════════════════════════════════════════════════════════
\section{Conclusions}
\label{sec:conclusions}

We identified a fundamental flaw in probability-weighted Monte Carlo
forecasters: even when paths are conditioned on historical patterns and
scored by fractal consistency (Hurst match, volatility match, smoothness),
the scoring criteria are symmetric with respect to direction.  The weighted
mean of such an ensemble discards the directional signal encoded in the
autocovariance structure.

The resolution is a split architecture that assigns each component its
optimal tool:
\begin{itemize}
  \item \textbf{Point forecast}: The Wiener--Hopf BLUP, solving
        $\boldsymbol{\Gamma}\mathbf{a} = \boldsymbol{\gamma}_h$ with
        adaptive lookback driven by autocovariance decay.
  \item \textbf{Prediction intervals}: Probability-weighted quantiles of
        pattern-conditioned MC paths scored by fractal consistency---distribution-free
        and well-calibrated.
\end{itemize}

On controlled synthetic data with known Hurst exponents, this split
forecaster (v3) achieves:
\begin{itemize}
  \item Lower RMSE than ARIMA at all tested $H$ (ratios 0.96--0.98).
  \item Higher directional accuracy at all $H$, consistently above 50\%.
  \item Higher Sharpe ratio at $H \geq 0.4$.
  \item Well-calibrated 95\% coverage (90.8--96.7\%).
\end{itemize}

The key contribution is not a new model but a \emph{structural insight}:
point prediction and uncertainty quantification should be solved separately
when the optimal tools for each are different.  The fGn autocovariance
provides the optimal linear predictor; the FracTime simulation pipeline---pattern-conditioned
paths, fractal consistency scoring, probability-weighted quantiles---provides
the optimal interval estimator.  Combining them yields a forecaster that is
simultaneously sharper and better calibrated than either ARIMA or the
original path-averaging approach.

% ══════════════════════════════════════════════════════════════════════════════
\bibliographystyle{plainnat}

\begin{thebibliography}{10}

\bibitem[Galbo(2026)]{Galbo2026}
Galbo, R. (2026).
\newblock The {FracTime} framework: A unified toolkit for fractal geometry and
  probability-weighted time series forecasting.
\newblock In \emph{AIRCC Conference Proceedings}.

\bibitem[Hurst(1951)]{Hurst1951}
Hurst, H.~E. (1951).
\newblock Long-term storage capacity of reservoirs.
\newblock \emph{Transactions of the American Society of Civil Engineers},
  116:770--808.

\bibitem[Mandelbrot and {Van Ness}(1968)]{MandelbrotVanNess1968}
Mandelbrot, B.~B. and {Van Ness}, J.~W. (1968).
\newblock Fractional {B}rownian motions, fractional noises and applications.
\newblock \emph{SIAM Review}, 10(4):422--437.

\bibitem[Peters(1994)]{Peters1994}
Peters, E.~E. (1994).
\newblock \emph{Fractal Market Analysis: Applying Chaos Theory to Investment
  and Economics}.
\newblock John Wiley \& Sons.

\bibitem[Box et~al.(2015)]{BoxJenkinsReinsel2015}
Box, G.~E.~P., Jenkins, G.~M., Reinsel, G.~C., and Ljung, G.~M. (2015).
\newblock \emph{Time Series Analysis: Forecasting and Control}.
\newblock John Wiley \& Sons, 5th edition.

\bibitem[Cont(2001)]{Cont2001}
Cont, R. (2001).
\newblock Empirical properties of asset returns: stylized facts and statistical
  issues.
\newblock \emph{Quantitative Finance}, 1(2):223--236.

\bibitem[Diebold and Mariano(1995)]{DieboldMariano1995}
Diebold, F.~X. and Mariano, R.~S. (1995).
\newblock Comparing predictive accuracy.
\newblock \emph{Journal of Business \& Economic Statistics}, 13(3):253--263.

\bibitem[Peng et~al.(1994)]{Peng1994}
Peng, C.-K., Buldyrev, S.~V., Havlin, S., Simons, M., Stanley, H.~E., and
  Goldberger, A.~L. (1994).
\newblock Mosaic organization of {DNA} nucleotides.
\newblock \emph{Physical Review E}, 49(2):1685--1689.

\bibitem[Gneiting and Raftery(2007)]{GneitingRaftery2007}
Gneiting, T. and Raftery, A.~E. (2007).
\newblock Strictly proper scoring rules, prediction, and estimation.
\newblock \emph{Journal of the American Statistical Association},
  102(477):359--378.

\bibitem[Kantelhardt et~al.(2002)]{Kantelhardt2002}
Kantelhardt, J.~W., Zschiegner, S.~A., Koscielny-Bunde, E., Havlin, S.,
  Bunde, A., and Stanley, H.~E. (2002).
\newblock Multifractal detrended fluctuation analysis of nonstationary time
  series.
\newblock \emph{Physica A}, 316(1--4):87--114.

\end{thebibliography}

\end{document}
